\documentclass[times]{elsarticle}
% \documentclass[review,doubleblind,times,12p]{elsarticle}
\let\today\relax
\makeatletter
\newcommand{\rmnum}[1]{\romannumeral #1}
\newcommand{\Rmnum}[1]{\expandafter\@slowromancap\romannumeral #1@}
\makeatother

\usepackage[margin=2.5cm]{geometry}
\usepackage{lineno}
\modulolinenumbers[5]
\usepackage{amsmath}
\usepackage{amsfonts,tabularx}
\usepackage{algorithm}
\usepackage[noend]{algpseudocode}

\usepackage{amsmath,amssymb,bm,graphicx,epstopdf}
\usepackage{amsthm}
\allowdisplaybreaks
\usepackage{dblfloatfix}
\usepackage{soul}
\usepackage[dvipsnames]{xcolor}
\usepackage{xurl}
\usepackage{xcolor}
\usepackage{amssymb}
\usepackage{fancyvrb}
\usepackage{soul}

\usepackage{booktabs,enumitem,ragged2e}

\usepackage{multirow}
\usepackage{graphicx}
\usepackage{subcaption}
\usepackage{graphbox}
\usepackage{mathtools}

%% Begin nomenclature setup

\usepackage{framed} % Framing content
\usepackage{multicol} % Multiple columns environment

\usepackage{nomencl} % Nomenclature package
\makenomenclature
\setlength{\nomitemsep}{-\parskip} % Baseline skip between items

% \renewcommand*\nompreamble{\begin{multicols}{2}}
% \renewcommand*\nompostamble{\end{multicols}}

\usepackage{etoolbox}
\renewcommand\nomgroup[1]{%
  \item[\bfseries
  \ifstrequal{#1}{A}{Abbreviations}{%
  \ifstrequal{#1}{I}{Indices and Constants}{}%
  \ifstrequal{#1}{V}{Variables and Functions}{}}
]}

%% End nomenclature setup
%\documentclass{article}
\usepackage{amsmath}
\newenvironment{conditions}
  {\par\vspace{\abovedisplayskip}\noindent\begin{tabular}{>{$}l<{$} @{${}={}$} l}}
  {\end{tabular}\par\vspace{\belowdisplayskip}}

\setlength{\parindent}{12pt}

\usepackage[noend]{algpseudocode}

\makeatletter
\newcommand{\multiline}[1]{%
  \begin{tabularx}{\dimexpr\linewidth-\ALG@thistlm}[t]{@{}X@{}}
    #1
  \end{tabularx}
}
\makeatother

\usepackage[colorlinks=true]{hyperref}
\hypersetup{
    colorlinks,
    citecolor=blue,
    linkcolor=blue
}
% \newif\ifwidefig

% \widefigfalse
%\widefigtrue



\newcommand{\eV}[2]{\ensuremath{\bar{V}_{#1}^{#2}}}


\newcommand{\red}[1]{{\color{red}#1}}

\DeclareMathOperator*{\argmax}{\arg\!\max}

% Additional packages/commands required by the cjc-2 manuscript content
\usepackage{mathrsfs}
\DeclareMathOperator*{\argmin}{\arg\!\min}

% Elsevier-compatible replacement for IEEEdescription used in cjc-2
\newenvironment{manuscriptdescription}
  {\begin{description}[style=multiline,leftmargin=4.8cm,labelwidth=4.5cm,align=left]}
  {\end{description}}

\begin{document}
\begin{frontmatter}

\title{Feasibility-Aware Scenario Learning for Carbon-Constrained Battery Planning in Data Centers}

\author[seuaddress]{Author Name\corref{correspondingauthor}}
\cortext[correspondingauthor]{Corresponding author.}
\ead{author@seu.edu.cn}
\address[seuaddress]{School of Electrical Engineering, Southeast University,
Nanjing, China}

\begin{abstract}
Battery planning for data centers must jointly account for electricity cost,
service continuity, and the carbon content of delivered electricity. Its
solution can be sensitive to rare combinations of high demand, low photovoltaic
generation, carbon-intensive grid supply, and grid outages, whereas conventional
scenario reduction is designed primarily to preserve statistical similarity.
This paper formulates a two-stage planning problem for a data center connected
through a single point of common coupling and introduces a source- and
charging-time-resolved ledger for stored carbon. The formulation tracks initial,
photovoltaic, diesel, and grid-charged energy layers and enforces hourly
carbon-intensity limits on electricity delivered to the data center. To
construct a compact scenario support, a conditional variational autoencoder is
fine-tuned using fixed-design recourse feedback from the exact planning model.
The proposed loss combines an infeasibility penalty, which corrects optimistic
demand, photovoltaic, and carbon trajectories, with an optimality-preserving
penalty that prevents overly conservative scenarios from excluding a known good
design. Mixed-integer solver outputs are used only as detached violation
severities; no differentiation through the solver or regret gradient is
required. The evaluation compares the learned support with statistical
scenario-selection and predict-then-optimize baselines in terms of
out-of-sample cost, load shedding, carbon excess, and storage design quality.
% TODO: Replace the final sentence with quantitative results.
\end{abstract}

\begin{keyword}
Battery energy storage \sep carbon accounting \sep data center \sep
decision-focused learning \sep scenario generation \sep stochastic planning
\end{keyword}

\end{frontmatter}

% =====================================================================
\section*{Nomenclature}

\subsubsection*{Sets and Indices}
\begin{manuscriptdescription}
\item[$\mathcal{S}$] Set of planning scenarios, indexed by $s$.
\item[$\mathcal{T}$] Set of hourly intervals, indexed by $t$.
\item[$\mathcal{K}_s$] Set of storage carbon layers, indexed by $k$.
\item[$\widehat{\mathcal{S}}$] Reduced scenario support.
\end{manuscriptdescription}

\subsubsection*{Parameters}
\begin{manuscriptdescription}
\item[$D_{s,t},V_{s,t}$] Data-center demand and available PV power.
\item[$a_{s,t}$] Grid-availability indicator.
\item[$p_{s,t},c_{s,t}$] Grid price and grid carbon intensity.
\item[$\pi_s,\rho_s$] Scenario probability and annual occurrence factor.
\item[$\eta^{\mathrm{ch}},\eta^{\mathrm{dc}}$] Storage charging/discharging efficiencies.
\item[$\underline\sigma,\overline\sigma$] Minimum/maximum state-of-charge fractions.
\item[$\bar c^{\mathrm n},\bar c^{\mathrm o}$] Normal/outage carbon limits.
\item[$\gamma_{s,k}$] Internal carbon intensity of storage layer $k$.
\end{manuscriptdescription}

\subsubsection*{Variables and Learning Symbols}
\begin{manuscriptdescription}
\item[$u^{\mathrm{es}}$] Binary storage build variable.
\item[$P^{\mathrm{cap}},E^{\mathrm{cap}}$] Installed power and energy capacity.
\item[$e_{s,k,t},P^{\mathrm{dc}}_{s,k,t}$] Layer energy and discharge power.
\item[$P^{\mathrm{shed}}_{s,t},P^{\mathrm{curt}}_{s,t}$] Load shedding and PV curtailment.
\item[$\Xi_{s,t}$] Carbon delivered above the hourly budget.
\item[$G_\theta$] Conditional scenario decoder.
\item[$L_{\mathrm{IPL}},L_{\mathrm{OPL}}$] Infeasibility and optimality-preserving losses.
\item[$\widehat{\bm y},\bm y^{\mathrm{PI}}$] Generated-support and perfect-information designs.
\item[$R$] Normalized out-of-sample decision regret.
\end{manuscriptdescription}

% =====================================================================
\section{Introduction}
\label{sec:introduction}

Electricity demand from data centers is large, persistent, and increasingly
exposed to both energy-system and environmental constraints. Behind-the-meter
photovoltaic (PV) generation and battery energy storage can reduce peak imports,
maintain service during grid interruptions, and shift electricity consumption
toward lower-carbon hours. These services interact. A battery sized only for
energy arbitrage may be inadequate during an outage, while a battery sized only
for reliability may be charged at times that undermine a carbon-intensity
target. Planning therefore requires a common representation of cost,
reliability, and carbon delivery under uncertainty.

Carbon-aware operation introduces an accounting question that is absent from
ordinary energy balancing. Electricity stored at different times may carry
different carbon intensities, and electricity charged from PV, the grid, and a
backup generator should not become indistinguishable merely because it occupies
the same physical battery. A perfectly mixed model assigns one average
intensity to the complete inventory. This abstraction can erase the temporal
and source information that gives storage its carbon value. Carbon-aware
optimal power-flow research has established the importance of attributing
emissions to electricity consumption \cite{chen2025copf}; an analogous issue
arises across time when electricity is carried through storage.

The second challenge is computational. Capacity planning is evaluated over
joint trajectories of demand, renewable availability, electricity price, grid
carbon intensity, and contingencies. A large historical pool provides coverage
but makes the mixed-integer planning problem expensive. Classical scenario
reduction chooses a smaller probability measure that is close to the original
under a distributional metric \cite{Dupacova2003}. Such proximity is not
necessarily aligned with the downstream decision. For carbon-constrained
storage, a modest statistical error can be decisive if it occurs during an
outage or in a charging hour whose energy is later used under a binding carbon
limit.

Decision-focused learning addresses the mismatch between predictive accuracy
and optimization quality \cite{Wilder2019}. Direct differentiation is difficult
here because the planning problem contains build and operating binaries and
because only the first-stage storage design, rather than a complete operating
schedule, transfers from generated scenarios to realized conditions. Moreover,
predicted demand, PV, and carbon intensity enter constraints. Optimistic
predictions may therefore produce a low-cost design that sheds load or violates
a carbon target when re-dispatched on real data. Feasibility-aware
decision-focused learning distinguishes such predictions from overly
conservative predictions that exclude a good decision \cite{Mandi2025}.

This paper develops an integrated planning and learning framework. A
conditional variational autoencoder (CVAE) \cite{Sohn2015} generates a compact
set of 48-hour scenarios. Each generated set is passed to an exact single-PCC
storage-planning model. Its build, power, and energy decisions are fixed and
re-evaluated with optimal recourse on real scenarios. Hourly load shedding and
carbon excess returned by the solver identify decision-critical prediction
errors. A complementary optimality-preserving term evaluates a
perfect-information design under the generated constraints, preventing the
generator from obtaining feasibility by making every scenario uniformly severe.

The main contributions are as follows.
\begin{enumerate}[leftmargin=*,label=\arabic*)]
\item A two-stage storage-planning formulation is developed for a single-PCC
data-center microgrid with electricity cost, demand charges, outage reliability,
and load-side carbon limits. An exact source- and charging-time-resolved ledger
preserves the carbon intensity of stored electricity without the bilinear mixing
identity required by an average-inventory model.
\item A feasibility-aware decision-focused scenario-generation method is
adapted to two-stage storage planning. It combines fixed-design recourse,
infeasibility and optimality-preserving losses, and decoder-gradient balancing,
while treating the mixed-integer solver as a detached black box.
\item A leakage-free evaluation protocol compares learned and statistically
selected supports on common held-out scenarios using cost, reliability, carbon
excess, installed capacity, and a perfect-information regret reference.
\end{enumerate}

Section~\ref{sec:formulation} formulates the storage-planning problem.
Section~\ref{sec:methodology} presents the decision-focused methodology.
Section~\ref{sec:case} describes the case study, and
Section~\ref{sec:conclusion} concludes the paper.

% =====================================================================
\section{Carbon-Aware Storage Planning Formulation}
\label{sec:formulation}

\subsection{System and Scenario Structure}

The system consists of a data center, behind-the-meter PV, a dispatchable
backup generator, and one candidate battery connected behind a single point of
common coupling (PCC). The PCC abstraction deliberately removes feeder
power-flow and multi-location siting effects. The first-stage decision is
whether to build the battery and, if built, its rated power and energy. All
dispatch variables are scenario dependent.

Scenario $s\in\mathcal S$ is the joint trajectory
\begin{equation}
\bm\xi_s=\{D_{s,t},V_{s,t},p_{s,t},c_{s,t},a_{s,t}\}_{t\in\mathcal T},
\label{eq:scenario}
\end{equation}
where $D$, $V$, $p$, $c$, and $a$ denote facility demand, available PV,
grid price, grid carbon intensity, and grid availability. Each block contains
$T=48$ hourly intervals. With physical probability $\pi_s$, its annual
occurrence factor is
\begin{equation}
\rho_s=\pi_s\frac{8760}{T\Delta t}.
\label{eq:annualization}
\end{equation}
This weighting corrects the deliberate oversampling of outage windows.

The planning problem has the two-stage form
\begin{equation}
\min_{\bm y\in\mathcal Y} I(\bm y)+
\mathbb{E}_{\bm\xi\sim\mathbb{P}}[Q(\bm y,\bm\xi)],
\qquad
\bm y=(u^{\mathrm{es}},P^{\mathrm{cap}},E^{\mathrm{cap}}).
\label{eq:two_stage}
\end{equation}

\subsection{Planning Objective}

The annualized investment cost is
\begin{equation}
I(\bm y)=\mathrm{CRF}
\left(c^{\mathrm{site}}u^{\mathrm{es}}+c^P P^{\mathrm{cap}}
+c^E E^{\mathrm{cap}}\right).
\label{eq:investment}
\end{equation}
Let $P^{\mathrm{peak}}$ denote the maximum grid import across planning
scenarios. The operating objective is
\begin{align}
C^{\mathrm{op}}={}&c^{\mathrm{dem}}P^{\mathrm{peak}}
+\sum_{s,t}\rho_s\Delta t\big[
p_{s,t}(P^{g,L}_{s,t}+P^{g,B}_{s,t})
\nonumber\\
&+c^d(P^{d,L}_{s,t}+P^{d,B}_{s,t})
+c^{\mathrm{deg}}(P^{\mathrm{ch}}_{s,t}+P^{\mathrm{dc}}_{s,t})
\nonumber\\
&+c^{\mathrm{curt}}P^{\mathrm{curt}}_{s,t}
+c^{\mathrm{shed}}P^{\mathrm{shed}}_{s,t}
+c^{\mathrm{CO_2}}\Xi_{s,t}\big].
\label{eq:operating_cost}
\end{align}
The planning oracle minimizes $C^{\mathrm{ann}}=I+C^{\mathrm{op}}$.
Carbon slack is diagnostic and penalized when enabled; otherwise the same
carbon condition is imposed as a hard constraint.

\subsection{Storage Design and PCC Dispatch}

The battery build, capacity, and duration limits are
\begin{align}
\underline P u^{\mathrm{es}}&\le P^{\mathrm{cap}}\le\overline P u^{\mathrm{es}},
\label{eq:power_cap}\\
\underline E u^{\mathrm{es}}&\le E^{\mathrm{cap}}\le\overline E u^{\mathrm{es}},
\label{eq:energy_cap}\\
h^{\min}P^{\mathrm{cap}}&\le E^{\mathrm{cap}}\le h^{\max}P^{\mathrm{cap}}.
\label{eq:duration}
\end{align}
Because there is one storage location, $u^{\mathrm{es}}$ is a build decision,
not a network siting decision.

Grid, PV, and diesel generation are split between direct load supply and
battery charging. PCC and PV balances are
\begin{align}
P^{g,L}_{s,t}+P^{v,L}_{s,t}+P^{d,L}_{s,t}+P^{\mathrm{dc}}_{s,t}
&=D_{s,t}-P^{\mathrm{shed}}_{s,t},
\label{eq:pcc_balance}\\
P^{v,L}_{s,t}+P^{v,B}_{s,t}+P^{\mathrm{curt}}_{s,t}
&=V_{s,t}.
\label{eq:pv_balance}
\end{align}
Resource and peak limits are
\begin{align}
P^{g,L}_{s,t}+P^{g,B}_{s,t}&\le a_{s,t}\overline P^g,
\label{eq:grid_limit}\\
P^{d,L}_{s,t}+P^{d,B}_{s,t}&\le\overline P^d,
\label{eq:diesel_limit}\\
P^{g,L}_{s,t}+P^{g,B}_{s,t}&\le P^{\mathrm{peak}}.
\label{eq:peak_limit}
\end{align}
Load shedding is restricted to outages:
\begin{equation}
0\le P^{\mathrm{shed}}_{s,t}\le D_{s,t}(1-a_{s,t}).
\label{eq:outage_shedding}
\end{equation}
Thus, the optimizer cannot meet a carbon limit by dropping load while connected.

Total charging and discharging are
\begin{align}
P^{\mathrm{ch}}_{s,t}&=P^{g,B}_{s,t}+P^{v,B}_{s,t}+P^{d,B}_{s,t},
\label{eq:total_charge}\\
P^{\mathrm{dc}}_{s,t}&=\sum_{k\in\mathcal K_s}P^{\mathrm{dc}}_{s,k,t}.
\label{eq:total_discharge}
\end{align}
With binary mode variables $z^{\mathrm{ch}}_{s,t}$ and
$z^{\mathrm{dc}}_{s,t}$,
\begin{align}
P^{\mathrm{ch}}_{s,t}&\le P^{\mathrm{cap}}, &
P^{\mathrm{ch}}_{s,t}&\le\overline Pz^{\mathrm{ch}}_{s,t},
\label{eq:charge_limit}\\
P^{\mathrm{dc}}_{s,t}&\le P^{\mathrm{cap}}, &
P^{\mathrm{dc}}_{s,t}&\le\overline Pz^{\mathrm{dc}}_{s,t},
\label{eq:discharge_limit}\\
z^{\mathrm{ch}}_{s,t}+z^{\mathrm{dc}}_{s,t}&\le u^{\mathrm{es}}.
\label{eq:mode}
\end{align}

\subsection{Source- and Vintage-Resolved Carbon Accounting}

The inventory is divided into virtual layers
\begin{equation}
\mathcal K_s=\{0,v,d\}\cup\{g_1,\ldots,g_T\},
\label{eq:layers}
\end{equation}
where $0$, $v$, and $d$ denote initial, PV-charged, and diesel-charged
inventory, while $g_k$ contains grid energy charged in interval $k$. These are
accounting layers rather than physically separated battery compartments.

Their fixed internal carbon intensities are
\begin{equation}
\gamma_{s,k}=
\begin{cases}
c^{\mathrm{init}},&k=0,\\
0,&k=v,\\
c^{d,\mathrm{CO_2}}/\eta^{\mathrm{ch}},&k=d,\\
c_{s,k}/\eta^{\mathrm{ch}},&k=g_k.
\end{cases}
\label{eq:layer_intensity}
\end{equation}
Charging losses reduce stored energy but do not erase source-attributed carbon,
which explains the division by $\eta^{\mathrm{ch}}$.

Let internal layer additions be
\begin{equation}
q_{s,k,t}=
\begin{cases}
\eta^{\mathrm{ch}}P^{v,B}_{s,t}\Delta t,&k=v,\\
\eta^{\mathrm{ch}}P^{d,B}_{s,t}\Delta t,&k=d,\\
\eta^{\mathrm{ch}}P^{g,B}_{s,t}\Delta t,&k=g_t,\\
0,&\text{otherwise}.
\end{cases}
\label{eq:additions}
\end{equation}
With retention $\kappa=1-\sigma^{\mathrm{sd}}$, layer energy satisfies
\begin{align}
e_{s,k,t+1}&=\kappa e_{s,k,t}+q_{s,k,t}
-\frac{P^{\mathrm{dc}}_{s,k,t}}{\eta^{\mathrm{dc}}}\Delta t,
\label{eq:layer_energy}\\
\frac{P^{\mathrm{dc}}_{s,k,t}}{\eta^{\mathrm{dc}}}\Delta t
&\le\kappa e_{s,k,t}.
\label{eq:available_energy}
\end{align}
Initial and aggregate SOC conditions are
\begin{align}
e_{s,0,1}&=\sigma^{\mathrm{init}}E^{\mathrm{cap}},
&e_{s,k,1}&=0,\quad k\ne0,
\label{eq:initial_layers}\\
\underline\sigma E^{\mathrm{cap}}
&\le\sum_ke_{s,k,t}\le\overline\sigma E^{\mathrm{cap}}.
\label{eq:soc}
\end{align}
The representative block returns to its initial energy and may not end with a
dirtier inventory:
\begin{align}
\sum_ke_{s,k,T+1}&=\sigma^{\mathrm{init}}E^{\mathrm{cap}},
\label{eq:terminal_energy}\\
\sum_k\gamma_{s,k}e_{s,k,T+1}
&\le c^{\mathrm{init}}\sigma^{\mathrm{init}}E^{\mathrm{cap}}.
\label{eq:terminal_carbon}
\end{align}

Carbon delivered to the data center is
\begin{equation}
C^L_{s,t}=c_{s,t}P^{g,L}_{s,t}
+c^{d,\mathrm{CO_2}}P^{d,L}_{s,t}
+\sum_k\gamma_{s,k}\frac{P^{\mathrm{dc}}_{s,k,t}}{\eta^{\mathrm{dc}}}.
\label{eq:delivered_carbon}
\end{equation}
Charging carbon enters the layer ledger and is assigned to load only when the
energy is discharged. This load-side compliance quantity is distinct from
source-side emissions, which also include energy used for charging.

\subsection{Carbon Limits and Benchmark Accounting}

The applicable intensity limit is
\begin{equation}
\bar c_{s,t}=a_{s,t}\bar c^{\mathrm n}+(1-a_{s,t})\bar c^{\mathrm o}.
\label{eq:applicable_cap}
\end{equation}
The main experiment imposes an hourly load-side budget:
\begin{equation}
C^L_{s,t}\le
\bar c_{s,t}(D_{s,t}-P^{\mathrm{shed}}_{s,t})+\Xi_{s,t},
\qquad \Xi_{s,t}\ge0.
\label{eq:hourly_cap}
\end{equation}
A horizon-average benchmark replaces it by
\begin{equation}
\sum_tC^L_{s,t}\Delta t\le
\sum_t\bar c_{s,t}(D_{s,t}-P^{\mathrm{shed}}_{s,t})\Delta t+\Xi_s.
\label{eq:horizon_cap}
\end{equation}

The layered formulation is linear apart from build and mode binaries. A
perfectly mixed benchmark tracks total energy $E_{s,t}$, carbon $M_{s,t}$, and
average intensity $w_{s,t}$ through
\begin{equation}
M_{s,t}=w_{s,t}E_{s,t},\qquad
C^B_{s,t}=w_{s,t}P^{\mathrm{dc}}_{s,t}/\eta^{\mathrm{dc}}.
\label{eq:average_model}
\end{equation}
These bilinear identities make the average-accounting benchmark nonconvex. It
is used only to test whether discarded source and vintage information changes
the reported design.

\subsection{Compact Formulation}

Let $\bm x_s$ collect all scenario-dependent variables and let
$\Omega(\bm\xi_s)$ denote
(\ref{eq:pcc_balance})--(\ref{eq:hourly_cap}). The deterministic equivalent is
\begin{equation}
\begin{aligned}
\min_{\bm y,\{\bm x_s\}}\quad
&I(\bm y)+\sum_{s\in\mathcal S}\rho_sC_s^{\mathrm{op}}(\bm x_s)\\
\mathrm{s.t.}\quad
&\bm y\in\mathcal Y,\\
&(\bm y,\bm x_s)\in\Omega(\bm\xi_s),\quad s\in\mathcal S.
\end{aligned}
\label{eq:deterministic_equivalent}
\end{equation}

% =====================================================================
\section{Decision-Focused Learning Methodology}
\label{sec:methodology}

\subsection{Learning Task and Scenario Generator}

Let $\mathcal D^{\mathrm{tr}}$ be the historical training pool and let
$K\ll|\mathcal D^{\mathrm{tr}}|$ be the number of scenarios affordable in the
planning model. The learned support is
\begin{equation}
\widehat{\mathcal S}_\theta=
\{\widehat{\bm\xi}_k=G_\theta(\bm z_k,\bm\omega_k)\}_{k=1}^K.
\label{eq:generated_support}
\end{equation}
Context $\bm\omega$ contains cyclical day-of-year coordinates, weekend share,
and mean temperature. Contexts and prior latent draws are fixed during
fine-tuning, so changes in the downstream design arise from generator updates
rather than resampling. The support also preserves the importance-corrected
normal/outage probability mass.

The CVAE encoder defines
\begin{equation}
q_\phi(\bm z\mid\bm x,\bm\omega)=
\mathcal N(\bm\mu_\phi,\operatorname{diag}(\bm\sigma_\phi^2)),
\label{eq:cvae}
\end{equation}
and its statistical objective is
\begin{align}
L_{\mathrm{CVAE}}={}&L_{\mathrm{rec}}+\lambda_rL_{\mathrm{ramp}}
+\lambda_nL_{\mathrm{net}}+\lambda_pL_{\mathrm{peak}}
\nonumber\\
&+\lambda_cL_{\mathrm{carbon\mbox{-}spread}}
+\beta D_{\mathrm{KL}}(q_\phi\Vert\mathcal N(\bm0,\bm I)).
\label{eq:cvae_loss}
\end{align}
Only demand, available PV, and grid carbon intensity are differentiable
decision channels. Workload and the calendar tariff are restored from legal
templates after decoding. Grid availability is copied as a detached exogenous
contingency, preventing the generator from manufacturing price arbitrage or
changing outage occurrence.

\subsection{Forward Planning and Fixed-Design Recourse}

The generated support induces the design
\begin{equation}
\widehat{\bm y}_\theta\in\argmin_{\bm y}
I(\bm y)+Q(\bm y;\widehat{\mathcal S}_\theta).
\label{eq:forward_plan}
\end{equation}
Only its build, power, and energy components are transferred to real scenario
$s$:
\begin{equation}
\bm x_s^R\in\argmin_{\bm x_s}
Q(\widehat{\bm y}_\theta,\bm x_s;\bm\xi_s).
\label{eq:fixed_recourse}
\end{equation}
All dispatch, battery-mode, layer-inventory, curtailment, shedding, and carbon
variables are reoptimized. Let $z^L_{s,t}$ and $z^C_{s,t}$ be the hourly
shedding and carbon excess returned by this solve. These quantities and all
MILP decisions are detached constants: they identify decision-critical errors
without defining a gradient through the optimizer.

\subsection{Infeasibility-Penalizing Loss}

For generated trajectories $\widehat D$, $\widehat V$, and $\widehat c$, define
\begin{equation}
d_{s,t}=\max(|D_{s,t}|,\varepsilon_D),\qquad
D^L_{s,t}=\max(D_{s,t}-z^L_{s,t},\varepsilon_D),
\label{eq:loss_scales}
\end{equation}
and optimistic load and PV directions
\begin{equation}
r^D_{s,t}=\frac{\left[D_{s,t}-\widehat D_{s,t}\right]_{+}}{d_{s,t}},\qquad
r^V_{s,t}=\frac{\left[\widehat V_{s,t}-V_{s,t}\right]_{+}}{d_{s,t}}.
\label{eq:load_pv_direction}
\end{equation}

Carbon is intertemporal because intensity in charging hour $k$ may affect
compliance in discharge hour $t$. Let $H^R_{s,t,k}$ be the detached exposure of
delivered carbon in hour $t$ to grid intensity in hour $k$, and define
\begin{equation}
\psi(c)=\tau\log\left(1+\exp((c-\bar c^{\mathrm n})/\tau)\right).
\label{eq:pressure}
\end{equation}
The exposure-weighted carbon direction is
\begin{equation}
r^C_{s,t}=
\frac{\left[\sum_kH^R_{s,t,k}[\psi(c_{s,k})-\psi(\widehat c_{s,k})]\right]_{+}}
{\bar c^{\mathrm n}\max(\sum_kH^R_{s,t,k},\varepsilon_H)}.
\label{eq:carbon_direction}
\end{equation}
This smooth pressure retains a correction near the cap, where a hard hinge may
otherwise be flat.

For $h_m(r)=\operatorname{softplus}(m+r)-\operatorname{softplus}(m)$,
\begin{align}
L_D^{\mathrm{IPL}}&=\frac1{KT}\sum_{s,t}
\frac{z^L_{s,t}}{d_{s,t}}h_m(r^D_{s,t}),
\label{eq:ipl_load}\\
L_V^{\mathrm{IPL}}&=\frac1{KT}\sum_{s,t}
\frac{z^L_{s,t}}{d_{s,t}}h_m(r^V_{s,t}),
\label{eq:ipl_pv}\\
L_C^{\mathrm{IPL}}&=\frac1{KT}\sum_{s,t}
\frac{z^C_{s,t}}{D^L_{s,t}}h_m(r^C_{s,t}),
\label{eq:ipl_carbon}
\end{align}
and
\begin{equation}
L_{\mathrm{IPL}}=w_DL_D^{\mathrm{IPL}}+w_VL_V^{\mathrm{IPL}}
+w_CL_C^{\mathrm{IPL}}.
\label{eq:ipl}
\end{equation}
Gradient descent therefore increases underestimated demand, decreases
overestimated PV, and increases underestimated carbon pressure only where the
generated-support design violates real recourse constraints.

\subsection{Optimality-Preserving Loss}

IPL alone can make every generated scenario severe. A perfect-information
design for the matched real scenario bag is therefore computed as
\begin{equation}
\bm y^{\mathrm{PI}}\in\argmin_{\bm y}
I(\bm y)+Q(\bm y;\mathcal S^R).
\label{eq:pi_design}
\end{equation}
This design is fixed and re-dispatched against the generated support. If the
generated constraints cause shedding or carbon excess, they have excluded a
design known to be adequate for the real bag. The optimality-preserving loss
uses reverse directions
\begin{equation}
\widetilde r^D_{s,t}=
\frac{\left[\widehat D_{s,t}-D_{s,t}\right]_{+}}{d_{s,t}},\qquad
\widetilde r^V_{s,t}=
\frac{\left[V_{s,t}-\widehat V_{s,t}\right]_{+}}{d_{s,t}},
\label{eq:opl_directions}
\end{equation}
and the exposure-weighted reverse of (\ref{eq:carbon_direction}). Substituting
the violation severities obtained by dispatching $\bm y^{\mathrm{PI}}$ defines
$L_{\mathrm{OPL}}$. It reduces excessive demand and carbon predictions and
increases excessively low PV predictions only where they reject the reference
design.

The combined decision loss is
\begin{equation}
L_{\mathrm{DFL}}=\alpha L_{\mathrm{IPL}}+(1-\alpha)L_{\mathrm{OPL}},
\qquad 0\le\alpha\le1.
\label{eq:dfl_loss}
\end{equation}
This construction adapts the feasibility/optimality logic of
\cite{Mandi2025} to two-stage recourse; it is not a derivative of the MILP value
function.

\subsection{Gradient Balancing and Training}

The total fine-tuning objective is
\begin{equation}
L_{\mathrm{total}}=w_{\mathrm{stat}}L_{\mathrm{CVAE}}
+\lambda_{\mathrm{eff}}L_{\mathrm{DFL}}.
\label{eq:total_loss}
\end{equation}
Let $g_{\mathrm{stat}}$ and $g_{\mathrm{DFL}}$ be decoder-gradient norms. The
detached effective weight is
\begin{equation}
\lambda_{\mathrm{eff}}=\lambda\operatorname{clip}\left(
r_g\frac{g_{\mathrm{stat}}}{\max(g_{\mathrm{DFL}},\varepsilon_g)},
b_{\min},b_{\max}\right).
\label{eq:gradient_balance}
\end{equation}
Only decoder gradients are compared because direct decision feedback reaches
the model through the decoder.

\begin{algorithm}[t]
\caption{Feasibility-aware CVAE fine-tuning}
\label{alg:training}
\begin{algorithmic}[1]
\Require Pretrained CVAE, training/validation pools, planning oracle, support size $K$
\State Fix support contexts, prior latents, contingencies, and weights
\For{each fine-tuning epoch}
  \State Decode $\widehat{\mathcal S}_\theta$ and solve (\ref{eq:forward_plan})
  \State Re-dispatch $\widehat{\bm y}_\theta$ on matched real scenarios
  \State Solve (\ref{eq:pi_design}) and dispatch it on generated scenarios
  \State Form $L_{\mathrm{IPL}}$, $L_{\mathrm{OPL}}$, and $L_{\mathrm{DFL}}$
  \State Evaluate $L_{\mathrm{CVAE}}$ on an independent reconstruction batch
  \State Balance decoder gradients and update $\theta$
  \State Evaluate and rank the fixed support on validation data
\EndFor
\State Restore the best validation checkpoint and return its $K$ scenarios
\end{algorithmic}
\end{algorithm}

Reliability is a checkpoint guardrail. Among checkpoints below a shedding
tolerance, lower carbon excess is preferred and exact decision regret breaks
ties. The pretrained model is a candidate at epoch $-1$, so fine-tuning is kept
only when it improves the fixed validation criterion.

\subsection{Out-of-Sample Metrics}

On common test support $\mathcal S^{\mathrm{te}}$, define
\begin{align}
C^{\mathrm{PI}}&=\min_{\bm y}I(\bm y)+Q(\bm y;\mathcal S^{\mathrm{te}}),
\label{eq:test_pi}\\
C_m&=I(\bm y_m)+Q(\bm y_m;\mathcal S^{\mathrm{te}}),
\label{eq:test_cost}\\
R_m&=\frac{C_m-C^{\mathrm{PI}}}{|C^{\mathrm{PI}}|+10^{-6}}.
\label{eq:regret}
\end{align}
Regret is detached and used only for evaluation. Storage value relative to the
same no-storage reference is $V_m=C^0-C_m$. Cost is reported together with
shedding, carbon excess, curtailment, and installed capacities.

% =====================================================================
\section{Case Study}
\label{sec:case}

\subsection{Data and Experimental Setup}

The study uses nonoverlapping 48-hour windows constructed from three years of
load, PV, temperature, tariff, and grid-carbon data. The active PCC reduction
retains the PV shape at the source feeder bus but synthesizes a 10-MW-class data
center trajectory from a workload proxy and temperature-dependent power-usage
effectiveness. Carbon trajectories remain intact 48-hour blocks and are
permuted only within season and split. Training, validation, and test sets are
disjoint.

Contiguous one- to four-hour outages are injected into a subset of windows so
small learning batches observe reliability events. Importance weights restore
an assumed physical outage probability of approximately $0.71\%$ per 48-hour
window. Outage timing is independent of load, PV, and carbon and is therefore a
stress-test assumption rather than a measured joint distribution.

\begin{table}[t]
\centering
\caption{Principal system and planning parameters}
\label{tab:parameters}
\begin{tabular}{@{}lr@{}}
\toprule
Parameter & Value \\
\midrule
Scenario horizon/resolution & 48 h/1 h \\
Mean/maximum demand & 8.477/10.224 MW \\
PV capacity & 5 MW \\
Grid-import limit & 15 MW \\
Backup-generator capacity & 6 MW \\
Battery power bound & 0--7.5 MW \\
Battery energy bound & 0--60 MWh \\
Initial/minimum/maximum SOC & 0.50/0.10/0.90 \\
Charge/discharge efficiency & 0.95/0.95 \\
Normal/outage carbon limit & 0.28/0.66 tCO$_2$/MWh \\
Planning support size & 8 \\
\bottomrule
\end{tabular}
\end{table}

The CVAE uses a 16-dimensional latent space and two 256-unit hidden layers in
each network. It is pretrained for 300 epochs and fine-tuned for 45 epochs. The
principal configuration uses eight fixed decision anchors, an independent
32-scenario reconstruction batch, $\alpha=0.5$, and decoder-gradient target
ratio $r_g=0.1$. Training planning uses a $1\%$ relative gap; final evaluation
uses the tighter oracle setting.

\subsection{Compared Methods and Metrics}

The comparison includes random, K-means, farthest-point, and aggregate observed
scenario supports; CVAE-only predict-then-optimize; IPL-only; OPL-only; and the
full DFL method. Direct optimization on the common test support is used only as
the perfect-information regret reference. All methods use the same support size
and their designs are re-dispatched on the same test scenarios.

\begin{table*}[t]
\centering
\caption{Out-of-sample planning performance on a common test set}
\label{tab:main_results}
\begin{tabular}{@{}lrrrrrrr@{}}
\toprule
Method & $P^{\mathrm{cap}}$ & $E^{\mathrm{cap}}$ & Annual cost & Regret
& Shedding & Carbon excess & Storage value \\
& (MW) & (MWh) & (\$) & (\%) & (MWh) & (tCO$_2$) & (\$) \\
\midrule
Random & -- & -- & -- & -- & -- & -- & -- \\
K-means & -- & -- & -- & -- & -- & -- & -- \\
Farthest & -- & -- & -- & -- & -- & -- & -- \\
Aggregate & -- & -- & -- & -- & -- & -- & -- \\
CVAE-only & -- & -- & -- & -- & -- & -- & -- \\
IPL-only & -- & -- & -- & -- & -- & -- & -- \\
OPL-only & -- & -- & -- & -- & -- & -- & -- \\
Full DFL & -- & -- & -- & -- & -- & -- & -- \\
Perfect information & -- & -- & -- & 0 & -- & -- & -- \\
\bottomrule
\end{tabular}
\end{table*}

% TODO: Populate from multi-seed evaluation artifacts. Report mean and standard
% deviation and distinguish incumbents from certified bounds where necessary.

The result discussion will answer whether decision-focused training changes the
installed design, whether any cost reduction is obtained at the expense of
physical violations, and whether improvements exceed both seed variability and
the uncertainty implied by solver gaps.

\subsection{Carbon-Accounting and Learning Ablations}

Four matched models compare layered/hourly, layered/horizon, average/hourly,
and average/horizon accounting. Identical data, weights, costs, and design
bounds must be used.

\begin{table}[t]
\centering
\caption{Effect of carbon accounting and cap scope}
\label{tab:carbon_ablation}
\begin{tabular}{@{}llrrr@{}}
\toprule
Accounting & Scope & $P^{\mathrm{cap}}$ & $E^{\mathrm{cap}}$ & Excess \\
& & (MW) & (MWh) & (tCO$_2$) \\
\midrule
Layered & Hourly & -- & -- & -- \\
Layered & Horizon & -- & -- & -- \\
Average & Hourly & -- & -- & -- \\
Average & Horizon & -- & -- & -- \\
\bottomrule
\end{tabular}
\end{table}

% TODO: Current t1--t4 YAML files are cost/DFL sensitivities, not the four
% accounting-scope combinations. Create matched configurations first.

The learning ablation varies $\alpha$, removes gradient balancing, and changes
support size. Planning sensitivities vary the normal carbon limit, battery
capital cost, and outage duration. Results should include means and standard
deviations over seeds. A representative 48-hour dispatch figure should show
grid intensity, delivered intensity, SOC, charging source, and discharge by
carbon layer to explain rather than merely report design differences.

\subsection{Computational Performance and Limitations}

Computational reporting includes forward-planning, fixed-design recourse, and
perfect-information times; solver calls per epoch; and final optimality gaps.
The current study cannot support conclusions about feeder congestion or
multi-location siting. Its demand is calibrated rather than metered, outage
timing is simulated independently of weather, and the directional DFL loss is
not the derivative of the MILP value function.

% =====================================================================
\section{Conclusion}
\label{sec:conclusion}

This paper formulated carbon-aware battery planning for a data center behind a
single PCC. The formulation jointly represents annualized cost, demand charges,
PV and diesel dispatch, outage-only load shedding, and load-side carbon limits.
Virtual inventory layers retain the source and charging-time intensity of
stored electricity while preserving the aggregate physical behavior of one
battery.

A feasibility-aware decision-focused method was developed to learn a compact
scenario support. It fixes the design obtained from generated scenarios and
evaluates it with optimal recourse on real scenarios. Infeasibility and
optimality-preserving losses use physical violations to correct optimistic and
overly conservative trajectories. Because all mixed-integer decisions and
severities are detached, the method avoids differentiating through the solver
while retaining a gradient path to generated demand, PV, and carbon.

% TODO: Insert quantitative findings after the result tables are populated.

Future work will extend the formulation to network-constrained multi-site
planning, correlated weather-driven outages, flexible computing workloads, and
longer horizons with inter-block storage states.

\section*{Acknowledgements}
% Add funding and non-author contributions.

\section*{CRediT authorship contribution statement}
Author Name: Conceptualization, Methodology, Software, Validation,
Writing--original draft.

\section*{Declaration of competing interest}
The authors declare that they have no known competing financial interests or
personal relationships that could have appeared to influence the work reported
in this paper.

\section*{Data availability}
The processed data, configurations, and source code used to reproduce the study
will be made available with the published article.

\bibliographystyle{elsarticle-num}
\bibliography{main_elsevier}
\end{document}
